\documentclass[10pt,a4paper]{amsart}
\usepackage[margin=19mm]{geometry}
\usepackage{amsmath,amssymb,mathtools}
\usepackage[scr=rsfs,cal=euler]{mathalfa}
\usepackage{xfrac}
\usepackage[unicode,hidelinks,bookmarks=false]{hyperref}
\newcommand{\ie}{i.e.\ }
\newcommand{\cf}{cf.\ }
\newcommand{\resp}{respectively\ }
\newcommand{\Subsection}{\subsection}
\providecommand{\nobreakdash}{\nobreak\hbox{-}}
\setlength{\emergencystretch}{2em}
\multlinegap0pt
\usepackage{amssymb}
\usepackage{mathtools}
\usepackage{bm}
\newtheorem{theorem}{Theorem}
\title{Bayesian Effective Dimension from Information Growth}
\author{Sayantan Banerjee}
\date{Source: arXiv:2512.23047v2 (CC BY 4.0)}
\begin{document}
\maketitle

\noindent\emph{Source and licence.} Bayesian Effective Dimension from Information Growth. Version arXiv:2512.23047v2 (CC BY 4.0), distributed by arXiv under the Creative Commons Attribution 4.0 International licence, http://creativecommons.org/licenses/by/4.0/, as recorded in the arXiv licence metadata for this exact version and shown on the abstract page https://arxiv.org/abs/2512.23047v2. Full licence notice: \texttt{licenses/arxiv-2512.23047-CC-BY-4.0.txt}.\par\smallskip

\noindent\textit{Editorial scope.} Counted unit: the article's theorem thm:exponential-spectrum (source0.tex lines 832--846). The article's complete original proof of it is delivered with it (source0.tex lines 848--937, one \texttt{proof} environment of 90 lines). No mathematics is written, repaired or re-derived: the statement and the proof are the article's own lines, in the article's own notation, and no retained line is substituted or re-flowed.  No further source-local context is needed: every cross-reference of every retained line resolves inside the two delivered spans.  Nothing was omitted from any retained span, so the package carries no omission record.  No retained line contains a citation, so the wrapper carries no bibliography and no bibliography entry.  No retained line contains a figure environment, an image-inclusion command or drawing code, so no image is delivered and the package declares no asset dependency.  Editorial formatting only, in addition: an AMS \texttt{amsart} preamble that restates the article's own declarations and macros used by the retained lines, namely the environments \texttt{theorem} on separate counters, together with the standard packages those lines need.  The retained environments print in this document as Theorem 1 in order of appearance, whereas the article prints the same items with the numbering of its own full document.  The complete native source of the article as distributed in the arXiv e-print for this exact version is bundled unchanged as \texttt{sources/191774/source0.tex}.
\begin{theorem}[Exponential spectral decay]
	\label{thm:exponential-spectrum}
	Suppose that \(-\log a_j\sim cj^\beta\), \(c>0\), \(\beta>0\). Then
	\[
	2I(\theta;Y)
	\sim
	\frac{\beta}{\beta+1}c^{-1/\beta}(\log n)^{1+1/\beta},
	\]
	and consequently
	\[
	d_{\rm eff}(n)
	\sim
	\frac{\beta}{\beta+1}c^{-1/\beta}(\log n)^{1/\beta}.
	\]
\end{theorem}

\begin{proof}
	We first consider the exact spectrum \(b_j=e^{-cj^\beta}\). Write \(L=\log n\) and \(m=(L/c)^{1/\beta}\). Then
	\[
	\sum_{j\geq1}\log(1+nb_j)
	=
	\sum_{j\geq1}\log\{1+\exp(L-cj^\beta)\}.
	\]
	
	Fix \(M>1\). For \(j\leq Mm\), \((L-cj^\beta)/L=1-(j/m)^\beta\). Using \(0\leq\log(1+e^z)-z_+\leq\log2\), \(z\in\mathbb R\), we obtain uniformly for \(j\leq Mm\),
	\[
	\frac{1}{L}\log\{1+\exp(L-cj^\beta)\}
	=
	\left\{1-\left(\frac{j}{m}\right)^\beta\right\}_+
	+
	O(L^{-1}).
	\]
	Hence
	\[
	\frac{1}{mL}
	\sum_{j\leq Mm}\log\{1+\exp(L-cj^\beta)\}
	\longrightarrow
	\int_0^M(1-x^\beta)_+\,dx
	=
	\frac{\beta}{\beta+1}.
	\]
	
	For the tail, \(\log(1+u)\leq u\) gives
	\[
	\sum_{j>Mm}\log\{1+\exp(L-cj^\beta)\}
	\leq
	e^L\sum_{j>Mm}e^{-cj^\beta}.
	\]
	Choose \(M_0\) with \(1<M_0<M\). For sufficiently large \(m\),
	\[
	\sum_{j>Mm}e^{-cj^\beta}
	\leq
	\sum_{j>M_0m}e^{-cj^\beta}
	\leq
	e^{-c\lfloor M_0m\rfloor^\beta}
	+
	\int_{\lfloor M_0m\rfloor}^{\infty}e^{-ct^\beta}\,dt.
	\]
	The standard incomplete-gamma estimate gives, for sufficiently large \(x\),
	\[
	\int_x^\infty e^{-ct^\beta}\,dt
	\leq
	C x^{1-\beta}e^{-cx^\beta}.
	\]
	Since \(c(M_0m)^\beta=M_0^\beta L\), multiplication by \(e^L\) yields
	\[
	e^L\sum_{j>Mm}e^{-cj^\beta}
	\leq
	\exp\{-(M_0^\beta-1)L+o(L)\}
	+
	C m^{1-\beta}
	\exp\{-(M_0^\beta-1)L+o(L)\}
	=
	o(mL),
	\]
	because \(M_0^\beta>1\). Thus, for the exact exponential spectrum,
	\[
	\sum_{j\geq1}\log(1+ne^{-cj^\beta})
	\sim
	\frac{\beta}{\beta+1}c^{-1/\beta}L^{1+1/\beta}.
	\]
	
	We now return to the assumed spectrum. For every \(\varepsilon\in(0,c)\), the relation \(-\log a_j\sim cj^\beta\) implies that, for all sufficiently large \(j\),
	\[
	e^{-(c+\varepsilon)j^\beta}
	\leq
	a_j
	\leq
	e^{-(c-\varepsilon)j^\beta}.
	\]
	The finitely many exceptional coordinates contribute \(O(\log n)=o\!\left((\log n)^{1+1/\beta}\right)\). Comparison with the two exact exponential spectra therefore gives
	\[
	\frac{\beta}{\beta+1}(c+\varepsilon)^{-1/\beta}
	\leq
	\liminf_{n\to\infty}
	\frac{2I(\theta;Y)}{(\log n)^{1+1/\beta}}
	\]
	and
	\[
	\limsup_{n\to\infty}
	\frac{2I(\theta;Y)}{(\log n)^{1+1/\beta}}
	\leq
	\frac{\beta}{\beta+1}(c-\varepsilon)^{-1/\beta}.
	\]
	Letting \(\varepsilon\downarrow0\) proves the first assertion. Division by \(\log n\) gives the result for \(d_{\rm eff}(n)\).
\end{proof}

\end{document}
